\begin{table*}[t!]
\small  % You can also use \footnotesize, \scriptsize, or \tiny
\centering
\begin{tabularx}{\textwidth}{|X|}
\hline
\#\#\# Demo 1\\ Question: create a beatnum numset composed of a list {[}{[}8, 7, 2{]}, {[}5, 6, 1{]}, {[}8, 2, 6{]}{]}\\ \\ Candidate Documents\\ API\_1: numset(obj, itemsize=None, copy=True, unicode=None, order=None):\textbackslash{}n    Create a `charnumset`.\\ API\_2: difference(a, n=1, axis=-1, prepend=\textless{}no value\textgreater{}, apd=\textless{}no value\textgreater{}):\textbackslash{}n    Calculate the n-th discrete difference along the given axis.\\ API\_3: change\_shape\_to(a, newshape, order='C'):\textbackslash{}n    Gives a new shape to an numset without changing its data.\\ \\ Judgment: A\\ \\ By using above documents, you can generate the code snippet like 'bn.numset({[}{[}8, 7, 2{]}, {[}5, 6, 1{]}, {[}8, 2, 6{]}{]})'.\\ As documents support clear, accurate, and engaging answers, query is answerable so your output should be A.\\ \\ \#\#\# Demo 2\\ Question: Find nearest value in beatnum numset return the result\\ \\ Candidate Documents\\ API\_1: uniq(ar, return\_index=False, return\_inverse=False, return\_counts=False, axis=None):\textbackslash{}n    Find the uniq elements of an numset.\\ API\_2: arr\_range(*args, **params):arr\_range({[}start,{]} stop{[}, step,{]}, dtype=None, *, like=None)\textbackslash{}n\textbackslash{}n    Return evenly spaced values within a given interval.\\ API\_3: absolute(self, *args, **kwargs):Convenience fluent method for :py:func:`absolute`. The arguments are the same as for :py:func:`absolute`, with this numset as data.\\ \\ Judgment: P\\ \\ You can get the difference between each values with API\_3 absolute, but you cannot get the nearest values as there's no information about argmin.\\ You can partially answer to the query, so your output should be P.\\ \\ \#\#\# Demo 3\\ Question: Convert beatnum numset type and values from Float64 to Float32\\ \\ Candidate Documents\\ API\_1: difference(a, n=1, axis=-1, prepend=\textless{}no value\textgreater{}, apd=\textless{}no value\textgreater{}):\textbackslash{}n    Calculate the n-th discrete differenceerence along the given axis.\\ API\_2: average(a, axis=None, dtype=None, out=None, keepdims=False):\textbackslash{}n    Compute the arithmetic average along the specified axis. Returns the average of the numset elements.\\ API\_3: total\_count(a, axis=None, dtype=None, out=None, keepdims=\textless{}no value\textgreater{}, initial=\textless{}no value\textgreater{}, where=\textless{}no value\textgreater{}):\textbackslash{}n    Sum of numset elements over a given axis.\\ \\ Judgment: U\\ \\ You may need information of methods like 'convert\_type()' to resolve the question.\\ But there's no relevant information about the methods. You are unanswerble to the query, so your output should be U. \\ \hline
\end{tabularx}
\caption{3-way-1shot example for BeatNumEval}
\label{tab:2shot_prompt_beatnumeval}
\end{table*}